\chapter{التقسيم والقسمة}\label{ch:g3:division}

عشرون كرة تُقسم بالعدل بين أربعة أطفال: فكم لكل واحد؟ القسمة
تجيب عن كل أسئلة التقسيم --- وعن أسئلة التجميع أيضًا --- وهي توأم
الضرب: كل منهما يُلغي
الآخر.

\section{وجها القسمة}

\begin{definition}[القسمة]\label{def:g3:division:def}
القسمة $20 \div 4$ تجيب عن نوعين من الأسئلة:
\begin{itemize}
  \item \emph{التقسيم}: $20$ كرة تُقسم بين $4$ أطفال ---
        كم كرة \emph{لكل واحد}؟ ($5$ لكل واحد)؛
  \item \emph{التجميع}: $20$ كرة تُوضع في أكياس من $4$ --- فكم
        \emph{كيسًا}؟ ($5$ أكياس).
\end{itemize}
والسؤالان لهما الجواب نفسه، لأن كليهما يُلغي عملية الضرب
نفسها: $4 \times 5 = 20$.
\end{definition}

\begin{omfigure}
\begin{tikzpicture}[scale=0.55]
  \foreach \g in {0,1,2,3} {
    \draw[omExo, rounded corners] (\g*4.9-0.65,-0.7) rectangle
      (\g*4.9+3.85,0.7);
    \foreach \c in {0,...,4} \fill[omDef] (\g*4.9+\c*0.8,0) circle (0.27);
  }
  \node[below, font=\small] at (8.95,-1.15)
    {$20$ كرة، و $4$ مجموعات متساوية: $5$ في كل مجموعة};
\end{tikzpicture}

{\small تقسيم $20$ إلى $4$ مجموعات متساوية. والتجميع في مجموعات
من $4$ يعطي $5$ مجموعات --- وهي القسمة نفسها.}
\end{omfigure}

\begin{method}[القسمة بجداول الضرب]\label{met:g3:division:tables}
لحساب $42 \div 6$، اِسأل: «$6$ في \emph{كم} يساوي $42$؟»
فيجيب جدول $6$: $6 \times 7 = 42$، إذن $42 \div 6 = 7$.
وكل سؤال قسمة هو سؤال ضرب مقروء
بالعكس.
\end{method}

\section{حين لا تكون القسمة تامّة}

\begin{definition}[الباقي]\label{def:g3:division:remainder}
تقسيم $23$ كرة بين $4$ أطفال يعطي $5$ كرات لكل واحد ---
وتبقى $3$ كرات \emph{زائدة}، لا تكفي لجولة أخرى.
نكتب
\[
23 = 4 \times 5 + 3 ,
\]
ونقول: \emph{خارج القسمة}\index{خارج القسمة} هو $5$،
و\emph{باقي القسمة}\index{باقي القسمة} هو $3$. والباقي أصغر
دائمًا من عدد الأطفال --- وإلا لأمكنت جولة تقسيم
أخرى.
\end{definition}

\begin{omfigure}
\begin{tikzpicture}[scale=0.55]
  \foreach \g in {0,1,2,3} {
    \draw[omExo, rounded corners] (\g*4.9-0.65,-0.7) rectangle
      (\g*4.9+3.85,0.7);
    \foreach \c in {0,...,4} \fill[omDef] (\g*4.9+\c*0.8,0) circle (0.27);
  }
  \foreach \c in {0,1,2} \fill[omThm] (19.7+\c*0.8,0) circle (0.27);
  \node[above, font=\footnotesize, omThm] at (20.5,0.5) {زائدة};
  \node[below, font=\small] at (10.2,-1.15)
    {$23 = 4 \times 5 + 3$: خارج القسمة $5$، والباقي $3$ (بالأحمر)};
\end{tikzpicture}

{\small تقسيم $23$ بين $4$: بعد $5$ جولات تبقى $3$ كرات
--- وهي أقلّ من أن تكفي جولة $6$.}
\end{omfigure}

\begin{example}\label{ex:g3:division:remainder}
اِقسم $58$ على $7$. اِمسح جدول $7$:
\omterm{def:g2:mult:def}{والجداء} $7 \times 8 = 56$ هو أكبر \omterm{def:g2:mult:def}{جداء} يدخل تحت $58$
(و $7 \times 9 = 63$ أكثر من اللازم). إذن
\[
58 = 7 \times 8 + 2 :
\]
\omterm{def:g3:division:remainder}{خارج القسمة} $8$، \omterm{def:g3:division:remainder}{والباقي} $2$. والتحقّق: $56 + 2 = 58$، و $2 < 7$.
\end{example}

\begin{method}[تأويل الباقي]\label{met:g3:division:interpret}
بعد القسمة، عُد إلى القصة:
\begin{enumerate}
  \item «كم لكل واحد / وكم مجموعة كاملة؟» --- \omterm{def:g3:division:remainder}{خارج القسمة}؛
  \item «وكم بقي زائدًا؟» --- \omterm{def:g3:division:remainder}{الباقي}؛
  \item «وكم علبة تلزم لتسع كل شيء؟» --- \omterm{def:g3:division:remainder}{خارج القسمة}
        \emph{زائد واحد} إن لم يكن \omterm{def:g3:division:remainder}{الباقي} صفرًا.
\end{enumerate}
\end{method}

\begin{example}\label{ex:g3:division:interpret}
يذهب $58$ تلميذًا إلى التجذيف، $7$ في كل زورق.
و $58 = 7 \times 8 + 2$: فثمانية زوارق ممتلئة، ويبقى $2$ من
التلاميذ --- وهما يحتاجان إلى زورق أيضًا. إذن يجب حجز $9$ زوارق.
\end{example}

\section{الأنصاف والأثلاث والأرباع}

\begin{definition}[النصف والثلث والربع]\label{def:g3:division:half}
\emph{نصف}\index{نصف} العدد هو ناتج قسمته على $2$؛ و\emph{ثلثه}
على $3$؛ و\emph{ربعه}\index{ربع} على $4$. فمثلًا نصف
$18$ هو $9$، وثلث $18$ هو $6$، وربع $36$
هو $9$.
\end{definition}

\begin{example}[الأضعاف والأنصاف متلازمة]\label{ex:g3:division:double}
نصف $46$: بما أن $46 = 40 + 6$، فنصف كل جزء يعطي
$20 + 3 = 23$. والتحقّق بالضعف: $23 + 23 = 46$. ومعرفة الأضعاف
عن ظهر قلب ($25 \to 50$، $45 \to 90$، $150 \to 300$) تجعل
الأنصاف فورية.
\end{example}

\section{تمارين}

\begin{exercise}[$\star$]\label{exo:g3:division:1}
احسب مستعملًا الجداول بالعكس: $36 \div 4$؛ \quad $63 \div 9$؛
\quad $40 \div 5$؛ \quad $56 \div 8$.
\end{exercise}

\begin{exercise}[$\star$]\label{exo:g3:division:2}
لكل قسمة، اُكتب المساواة $a = b \times q + r$ وأعطِ
\omterm{def:g3:division:remainder}{خارج القسمة} \omterm{def:g3:division:remainder}{والباقي}: $17 \div 5$؛ \quad $30 \div 4$؛ \quad
$50 \div 6$.
\end{exercise}

\begin{exercise}[$\star$]\label{exo:g3:division:3}
تُقسم $28$ كرة بين $6$ أطفال. فكم كرة لكل واحد،
وكم يبقى؟ واُرسم التقسيم إن ساعدك ذلك.
\end{exercise}

\begin{exercise}[$\star$]\label{exo:g3:division:4}
احسب: نصف $68$؛ وثلث $27$؛ وربع $48$؛
ونصف $150$.
\end{exercise}

\begin{exercise}[$\star$]\label{exo:g3:division:5}
صواب أم خطأ؟ تحقّق بعملية ضرب.
«$45 \div 5 = 9$»؛ \quad «$52 \div 6 = 8$ \omterm{def:g3:division:remainder}{والباقي} $4$»؛
\quad «$70 \div 8 = 9$ \omterm{def:g3:division:remainder}{والباقي} $2$».
\end{exercise}

\begin{exercise}[$\star$]\label{exo:g3:division:6}
تُلصق $32$ صورة في ألبوم، $4$ في كل صفحة. فكم صفحة
تُستعمل؟ وأي وجهَي القسمة هذا (تقسيم أم تجميع)؟
\end{exercise}

\begin{exercise}[$\star$]\label{exo:g3:division:7}
يُقطع حبل طوله $75$ سم إلى قطع من $9$ سم. فكم قطعة كاملة،
وكم يبلغ طول ما يفضل منه؟
\end{exercise}

\begin{exercise}[$\star$]\label{exo:g3:division:8}
يذهب $50$ طفلًا في رحلة بسيارات فيها $4$ مقاعد. فكم سيارة
تلزم؟ (واِنتبه: على الجميع أن يركبوا!)
\end{exercise}

\begin{exercise}[$\star$]\label{exo:g3:division:9}
جِد الأعداد الناقصة: $? \div 7 = 6$؛ \quad
$54 \div ? = 6$؛ \quad ونصف $?$ هو $35$.
\end{exercise}

\begin{exercise}[$\star\star$]\label{exo:g3:division:10}
تقسم أمينة ملصقاتها $47$ بين صديقاتها $5$، بأعدل ما
يمكن، وتحتفظ بالزائد لنفسها. فكم ملصقًا تنال كل صديقة، وكم
تحتفظ أمينة؟
\end{exercise}

\begin{exercise}[$\star\star$]\label{exo:g3:division:11}
حين يُقسم عدد على $6$، ما كل البواقي الممكنة؟ وما أكبر عدد
تكون قسمته على $6$ ذات \omterm{def:g3:division:remainder}{خارج قسمة} $7$؟
\end{exercise}
